\section{Tensor Product}

We've never actually defined so far, what a tensor is in mathematical terms, so we're going to do this now in the most correct
and abstract way: \textbf{A tensor is a collection of vectors and covectors combined together using the tensor product}.\\

This definition requires an explanation and understanding of what this tensor product is, as this is the first time 
we encounter the term.\\

We've seen what vectors and covectors are, how they transform when changing coordinate systems.
We've then see what a linear map is, and how that transform as well, and finally what a bilinear map is and again, how the components
transform.
What we're going to do now is level up the notions we've learned, and actually make the transformation rules comes out in a more 
elegant and easier way, knowing just the fundamental transform rules that basis vectors follow (forward and backward transformation), the rest
(linear maps, bilinear maps and eventually multi-linear maps) will come naturally from there, once we describe what the tensor product is.\\

If we think about tensors, you already know that the type is given by the contravariant and covariant components of it, and you know that
vectors are contravariant tensors (1,0) and covectors are covariant tensors (0,1) so you may already understand where we're going here: we'd like to 
be able to build any type of tensors using these basic building blocks (vectors and covectors) combined together in some form.

\subsection{Linear maps as vector-covector pairs}

The first tensor we're going to try to build is a linear map.
We know that a linear map is a (1,1)-tensor, that takes a vector space to another: $\textit{L} : V \mapsto V$ for example to keep it simple.

Now, linear maps are represented in coordinates as a matrix, so we may try thinking how we can obtain a matrix from vectors and covectors.
Well you know that by following the row-to-column matrix multiplication rule, you can write down a matrix as a column vector multiplied by a row vector.\\

You can imagine a simple example matrix 

\begin{equation}
    \begin{bmatrix}
        4 & 400\\
        8 & 800
    \end{bmatrix}
\end{equation}

as a representation of a linear map in a certain basis. And you can ask yourself, which are the column and row vectors I can use to build such a matrix?
If you try yourself, you will find out that this is possible by multiplying:

\begin{equation}
    \begin{bmatrix}
        4\\
        8
    \end{bmatrix}
    \begin{bmatrix}
        1 & 100
    \end{bmatrix} = \begin{bmatrix}
        4 & 400\\
        8 & 800
    \end{bmatrix}
\end{equation}

So there are matrices (linear maps) that can be broken up into a column vector and a row vector in this way: these are usually called pure linear maps.
Whereas there are others, that cannot be obtained by such a combination, these are called inpure.
Pure matrices are actually quite boring, because such linear maps applied to vectors, keep the output in the same direction.\\

Impure matrices instead, are the more interesting ones, because we can send input vectors into any direction, and we're not as limited as we're with pure ones.
So we have a sort of a problem here, it seems that we're able to only produce pure matrices from colum/row vector products, and they're the boring ones.\\
How can we construct inpure matrices from colum/row vector products then? \\

We can try defining four special vector-covector pairs, using the basis vectors and basis covectors (dual basis):

\begin{align*}
    \color{blue}{\vec{e_1}, \vec{e_2}} && \color{blue}{\epsilon^1, \epsilon^2}
\end{align*}

\begin{align*}
    \color{blue} \vec{e_1} \epsilon^1 &= 
    \begin{bmatrix}
        1\\
        0
    \end{bmatrix}
    \begin{bmatrix}
        1 & 0
    \end{bmatrix} = 
    \begin{bmatrix}
        1 & 0\\
        0 & 0
    \end{bmatrix}\\
    \color{blue} \vec{e_1} \epsilon^2 &= 
    \begin{bmatrix}
        1\\
        0
    \end{bmatrix}
    \begin{bmatrix}
        0 & 1
    \end{bmatrix} = 
    \begin{bmatrix}
        0 & 1\\
        0 & 0
    \end{bmatrix}\\
    \color{blue} \vec{e_2} \epsilon^1 &= 
    \begin{bmatrix}
        0\\
        1
    \end{bmatrix}
    \begin{bmatrix}
        1 & 0
    \end{bmatrix} = 
    \begin{bmatrix}
        0 & 0\\
        1 & 0
    \end{bmatrix}\\
    \color{blue} \vec{e_2} \epsilon^2 &= 
    \begin{bmatrix}
        0\\
        1
    \end{bmatrix}
    \begin{bmatrix}
        0 & 1
    \end{bmatrix} = 
    \begin{bmatrix}
        0 & 0\\
        0 & 1
    \end{bmatrix}
\end{align*}

With these basis vector-covector pairs above, we can basically write any linear maps (using indeed this basis) as a linear
combination:

\begin{equation}
    \textit{L} : \begin{bmatrix}
        a & b\\
        c & d
    \end{bmatrix} = 
    a
    \begin{bmatrix}
        1 & 0\\
        0 & 0
    \end{bmatrix} + b
    \begin{bmatrix}
        0 & 1\\
        0 & 0
    \end{bmatrix} + c
    \begin{bmatrix}
        0 & 0\\
        1 & 0
    \end{bmatrix} + d
    \begin{bmatrix}
        0 & 0\\
        0 & 1
    \end{bmatrix}
\end{equation}

So these four vector-covector pairs form a basis for any linear maps from the vector space $V$ to itself.
In Einstein notation, we can write:

\begin{equation}
    \textit{L} = \textit{L}^i_j \color{blue} \vec{e_i} \epsilon^j
\end{equation}

If we try to apply this linear map to a generic vector to obtain another one in the same space:

\begin{align*}
    \begingroup\color{violet}w\endgroup = \textit{L}\left(\begingroup\color{orange}\vec{v}\endgroup\right) &= 
    \textit{L}^i_j \color{blue} \vec{e_i} \epsilon^j \left(\begingroup\color{orange}v^k\endgroup \begingroup\color{blue}\vec{e_k}\endgroup\right)\\
    &= \textit{L}^i_j \begingroup\color{orange}v^k\endgroup \begingroup\color{blue}\vec{e_i}\endgroup \begingroup\color{blue}\epsilon^j\endgroup \left(\begingroup\color{blue}\vec{e_k}\endgroup\right)\\
    &= \textit{L}^i_j \begingroup\color{orange}v^k\endgroup \begingroup\color{blue}\vec{e_i}\endgroup \begingroup\color{blue}\epsilon^j\endgroup \delta^j_k\\
    &= \textit{L}^i_j \begingroup\color{orange}v^j\endgroup\begingroup\color{blue}\vec{e_i}\endgroup
\end{align*}

As you can see, that's out output vector written as linear combination of the basis vectors.
This shows that indeed, this vector-covector pair is a linear map, that takes a vector as input, and gets a vector as output.\\

We can say that these are a basis for $V \mapsto V$:

\begin{equation}
    {\color{blue} \vec{e_i}\epsilon^j = \vec{e_1} \epsilon^1,  \vec{e_1} \epsilon^2,  \vec{e_2} \epsilon^1, \vec{e_2} \epsilon^2}
\end{equation}

We chose this basis but we could have chosen a different one of course, this is just convenient in terms of calculation.\\
Also note that the notation used is not really standard, as usually you actually see the tensor product 
written correctly as:

\begin{equation}
    \begingroup\color{blue} \vec{e_i}\endgroup \otimes \begingroup\color{blue}\epsilon^j \endgroup
\end{equation}

And the linear map can be written like:
\begin{equation}
    \textit{L} = \textit{L}^i_j \begingroup\color{blue} \vec{e_i} \epsilon^j\endgroup = \textit{L}^i_j \left(\begingroup\color{blue} \vec{e_i}\endgroup \otimes \begingroup\color{blue}\epsilon^j \endgroup\right)
\end{equation}



\subsection{Bilinear forms as covector-covector pairs}

We're going to have a look at a similar perspective we used for linear maps, but applied to Bilinear forms now, metric tensor included.
Bilinear forms can be indeed thought as linear combination of covector-covector pairs.

\begin{equation}
    \mathcal{B} = \mathcal{B}_{ij} \begingroup\color{blue} \epsilon^i \epsilon^j\endgroup = \mathcal{B}_{ij} \left(\begingroup\color{blue} \epsilon^i\endgroup \otimes \begingroup\color{blue}\epsilon^j \endgroup\right)
\end{equation}

Why covector-covector pairs and not something else? Well the logic is that, a bilinear forms take two vector inputs, right?\\
Since covectors are basically linear maps that take each a vector input and outputs a scalar, a pair of covectors
would take two vectors as input.\\

Trying that out with the transformation rules, we'll see that this makes sense and we're going to end up
with the correct transformation we previously found for bilinear forms:

\begin{align*}
    \mathcal{B} &= \begingroup\color{cyan} B_{kl}\endgroup \begingroup\color{blue}\epsilon^k \epsilon^l\endgroup \\
    &= \begingroup\color{cyan} B_{kl}\endgroup \left(F^k_i \begingroup\color{red} \tilde{\epsilon^i}\endgroup\right) \left(F^l_j \begingroup\color{red} \tilde{\epsilon^j}\endgroup\right)\\
    &= \left( F^k_l F^l_j \begingroup\color{cyan} B_{kl}\endgroup \right) \begingroup\color{red} \tilde{\epsilon^i}\tilde{\epsilon^j}\endgroup\\
    &= \begingroup\color{magenta} \tilde{\mathcal{B}_{ij}} \endgroup \begingroup\color{red} \tilde{\epsilon^i}\tilde{\epsilon^j}\endgroup
\end{align*}

So the components transform as:

\begin{equation}
    \begingroup\color{magenta} \tilde{\mathcal{B}_{ij}} \endgroup =  F^k_l F^l_j \begingroup\color{cyan} B_{kl}\endgroup
\end{equation}

A similar process for the reverse transformation yields:

\begin{equation}
    \begingroup\color{cyan} B_{ij}\endgroup = B^k_i B^l_j \begingroup\color{magenta} \tilde{\mathcal{B}_{kl}} \endgroup 
\end{equation}

Next, we can also check that we get the correct component multiplication formula when we get this form applied to two vectors as input:

\begin{align*}
    \mathcal{B} &= B_{ij} \begingroup\color{blue} \epsilon^i \epsilon^j \endgroup\\
    \begingroup\color{orange}\vec{v}\endgroup &= \begingroup\color{orange} v^k\endgroup \begingroup\color{blue} \vec{e_k}\endgroup\\
    \begingroup\color{violet}\vec{w}\endgroup &= \begingroup\color{violet} w^l\endgroup \begingroup\color{blue} \vec{e_l}\endgroup\\ \\
    s &= \mathcal{B} \left(\begingroup\color{orange}\vec{v}\endgroup, \begingroup\color{violet}\vec{w}\endgroup\right) = B_{ij} \begingroup\color{blue}\epsilon^i \epsilon^j \endgroup \left(\begingroup\color{orange} v^k\endgroup \begingroup\color{blue} \vec{e_k}\endgroup, \begingroup\color{violet} w^l\endgroup \begingroup\color{blue} \vec{e_l}\endgroup\right)\\
    &= B_{ij} \begingroup\color{blue}\epsilon^i \endgroup \left(\begingroup\color{orange} v^k\endgroup \begingroup\color{blue} \vec{e_k}\endgroup\right) \begingroup\color{blue}\epsilon^j \endgroup \left(\begingroup\color{violet} w^l\endgroup \begingroup\color{blue} \vec{e_l}\endgroup\right)\\
    &= B_{ij} \begingroup\color{orange} v^k\endgroup \begingroup\color{violet} w^l\endgroup \delta^i_k \delta^j_l\\
    &= B_{ij} \begingroup\color{orange} v^i\endgroup \begingroup\color{violet} w^j\endgroup
\end{align*}

And this, if you remember, is the correct component multiplication formula that gives us the correct final single number as output.
